% concatenated from main.tex
\documentclass[11pt]{amsart}
\usepackage{tikz}
\usepackage{mathrsfs}
\usepackage{amscd,verbatim}
\usepackage{floatflt}
\usepackage{amssymb}
\usepackage{tikz-cd}
\usepackage[all]{xy}
%\usepackage{geometry}

\usepackage{stmaryrd} % for \mapsfrom
%\usepackage[color,notcite,notref]{showkeys}
%\definecolor{labelkey}{rgb}{1,0,0}
\usepackage[colorlinks,linkcolor=blue,citecolor=blue,urlcolor=red]{hyperref}
\hbadness=10000
\vbadness=10000
\hfuzz=100pt
%\renewcommand{\baselinestretch}{1.14}
\newcommand{\rem}[1]{{\color{red}\sout{#1}}}
\newcommand{\adj}[1]{{\color{blue}#1}}

\newcommand{\ulMCorf}{{\ulMCor}^{\text{fin}}}
\newcommand{\ulMMO}{\textbf{\underline{M}}\sO}
\newcommand{\ulMO}{\ul{M}\sO}
\newcommand{\MpU}{(\ol{U},U^{\infty})}
\newcommand{\un}{\mathbf{1}}
\newcommand{\A}{\mathbf{A}}
\renewcommand{\AA}{\mathbb{A}}
\newcommand{\mbA}{\mathbb{A}^1}
\newcommand{\mbAm}{\mathbb{A}^m}
\newcommand{\deAmi}{(\mathbb{A}^{1},d_{i_{1}}\{0\}) \otimes (\mathbb{A}^{1},d_{i_{2}}\{0\}) \otimes \cdots \otimes (\mathbb{A}^{1},\emptyset)}
\newcommand{\deAmoi}{(\mathbb{A}^{1},d_{i_{1}}\{0\}) \otimes (\mathbb{A}^{1},d_{i_{2}}\{0\}) \otimes \cdots \otimes \{0\}}
\newcommand{\deAm}{(\mathbb{A}^{1},d_{{1}}\{0\}) \otimes (\mathbb{A}^{1},d_{{2}}\{0\}) \otimes \cdots \otimes (\mathbb{A}^{1},\emptyset)}
\newcommand{\deAmo}{(\mathbb{A}^{1},d_{{1}}\{0\}) \otimes (\mathbb{A}^{1},d_{{2}}\{0\}) \otimes \cdots \otimes \{0\}}
\newcommand{\deAl}{(\mathbb{A}^{1},d_{1}\{0\})\otimes (\mathbb{A}^{1},d_{2}\{0\})\otimes \cdots \otimes (\mathbb{A}^{1},d_{m-1}\{0\})}
\newcommand{\BB}{\mathbb{B}}
\newcommand{\F}{\mathbf{F}}
\newcommand{\G}{\mathbf{G}}
\renewcommand{\P}{\mathbf{P}}
\renewcommand{\L}{\mathbb{L}}
\newcommand{\N}{\mathbb{N}}
\newcommand{\PP}{\mathbb{P}}
\newcommand{\Q}{\mathbb{Q}}
\newcommand{\Z}{\mathbb{Z}}
\newcommand{\sA}{\mathcal{A}}
\newcommand{\sB}{\mathcal{B}}
\newcommand{\sC}{\mathcal{C}}
\newcommand{\sD}{\mathcal{D}}
\newcommand{\sE}{\mathcal{E}}
\newcommand{\sF}{\mathcal{F}}
\newcommand{\sG}{\mathcal{G}}
\newcommand{\sH}{\mathcal{H}}
\newcommand{\sI}{\mathcal{I}}
\newcommand{\sJ}{\mathcal{J}}
\newcommand{\sK}{\mathcal{K}}
\newcommand{\sL}{\mathcal{L}}
\newcommand{\sO}{\mathcal{O}}
\newcommand{\sP}{\mathcal{P}}
\newcommand{\sQ}{\mathcal{Q}}
\newcommand{\sR}{\mathcal{R}}
\newcommand{\sS}{\mathcal{S}}
\newcommand{\sT}{\mathcal{T}}
\newcommand{\sU}{\mathcal{U}}
\newcommand{\sV}{\mathcal{V}}

\newcommand{\Zt}{\Ztr}
\newcommand{\Shv}{{\operatorname{Shv}}}
\newcommand{\sX}{\mathcal{X}}
\newcommand{\sY}{\mathcal{Y}}
\newcommand{\sZ}{\mathcal{Z}}
\newcommand{\bH}{\mathbb{H}}
\newcommand{\bN}{\mathbb{N}}
\newcommand{\bZ}{\mathbb{Z}}
\newcommand{\fm}{\mathfrak{m}}
\newcommand{\fM}{\mathfrak{M}}
\newcommand{\fU}{\mathfrak{U}}
\newcommand{\fV}{\mathfrak{V}}
\newcommand{\fW}{\mathfrak{W}}
\newcommand{\fX}{\mathfrak{X}}
\newcommand{\fY}{\mathfrak{Y}}
\newcommand{\fZ}{\mathfrak{Z}}
\newcommand{\fn}{\mathfrak{n}}
\newcommand{\fp}{\mathfrak{p}}
\newcommand{\fq}{\mathfrak{q}}
\newcommand{\fr}{\mathfrak{r}}
%\newcommand{\Cb}{{\overline{C}}}
\newcommand*{\xdashrightarrow}[2][]{%
  \mathrel{%
    \mathpalette{\da@xarrow{#1}{#2}{}\da@rightarrow{\,}{}}{}%
  }%
}
\newcommand{\Xb}{{\overline{X}}}
\newcommand{\Tb}{\overline{T}}
\newcommand{\alphab}{\ol{\alpha}}
\newcommand{\deltab}{\ol{\delta}}
\newcommand{\tilO}{\tilde{\mathcal{O}}}
\newcommand{\Mfget}{\ulM^{\text{h}}_{\text{gfin}}}
\newcommand{\Mhgf}{\ulM^{\text{h}}_{\text{gfin}}}
\newcommand{\Mqfh}{\ulM \text{qfh}}


\newcommand{\tilOmega}{\tilde{\Omega}}
\newcommand{\Ga}{{G^{(0)}}}
\newcommand{\Gb}{{G^{(1)}}}
\newcommand{\CuXY}{\sC_{(\Xb,Y)}}
\newcommand{\ha}{{h^{(0)}}}
\newcommand{\hb}{{h^{(1)}}}
\newcommand{\Mod}{\operatorname{Mod}\hbox{--}}
\newcommand{\Span}{\operatorname{\mathbf{Span}}}
\newcommand{\Cor}{\operatorname{\mathbf{Cor}}}
\newcommand{\Mack}{\operatorname{\mathbf{Mack}}}
\newcommand{\HI}{\operatorname{\mathbf{HI}}}
\newcommand{\Pre}{{\operatorname{\mathbf{Pre}}}}
\newcommand{\Rec}{{\operatorname{\mathbf{Rec}}}}
\newcommand{\SHI}{\operatorname{\mathbf{SHI}}}
\newcommand{\proj}{\operatorname{proj}}


\newcommand{\Cone}{\operatorname{Cone}}
\newcommand{\Ext}{\operatorname{Ext}}
\newcommand{\ul}[1]{{\underline{#1}}}

\newcommand{\Set}{{\operatorname{\mathbf{Set}}}}
\newcommand{\PST}{{\operatorname{\mathbf{PST}}}}
\newcommand{\PS}{{\operatorname{\mathbf{PS}}}}
\newcommand{\NST}{\operatorname{\mathbf{NST}}}
\newcommand{\RS}{\operatorname{\mathbf{RS}}}
\newcommand{\Chow}{\operatorname{\mathbf{Chow}}}
\newcommand{\DM}{\operatorname{\mathbf{DM}}}
\newcommand{\DMR}{\operatorname{\mathbf{DMR}}}
\newcommand{\DR}{\operatorname{\mathbf{DR}}}
\newcommand{\MDM}{\operatorname{\mathbf{MDM}}}
\newcommand{\MDMgm}{\MDM_{\text{gm}}^{\text{eff}}}
\newcommand{\ulMDM}{\operatorname{\mathbf{\underline{M}DM}}}
%\newcommand{\DR}{\operatorname{\mathbf{\underline{M}DM}}}
\newcommand{\Map}{\operatorname{Map}}
\newcommand{\Hom}{\operatorname{Hom}}
\newcommand{\uHom}{\operatorname{\underline{Hom}}}
\newcommand{\uExt}{\operatorname{\underline{Ext}}}
\newcommand{\RHom}{\operatorname{R\underline{Hom}}}
\newcommand{\que}[1]{{\color{green}#1}}
\newcommand{\uG}{{\underline{G}}}
\newcommand{\Ker}{\operatorname{Ker}}
\newcommand{\IM}{\operatorname{Im}}
\renewcommand{\Im}{\operatorname{Im}}
\newcommand{\Coker}{\operatorname{Coker}}
\newcommand{\MpN}{(\ol{N},N^{\infty})}
\newcommand{\Coim}{\operatorname{Coim}}
\newcommand{\Tr}{\operatorname{Tr}}
\newcommand{\Div}{\operatorname{Div}}
\newcommand{\Pic}{\operatorname{Pic}}
\newcommand{\Br}{\operatorname{Br}}
\newcommand{\Spec}{\operatorname{Spec}}
\newcommand{\Proj}{\operatorname{Proj}}
\newcommand{\Sm}{\operatorname{\mathbf{Sm}}}
\newcommand{\Sch}{\operatorname{\mathbf{Sch}}}
\newcommand{\Ab}{\operatorname{\mathbf{Ab}}}
\newcommand{\oo}{\operatornamewithlimits{\otimes}\limits}
\newcommand{\by}{\xrightarrow}
\newcommand{\yb}{\xleftarrow}
\newcommand{\iso}{\by{\sim}}
\newcommand{\osi}{\yb{\sim}}
\newcommand{\rec}{{\operatorname{rec}}}
\newcommand{\pro}[1]{\text{\rm pro}_{#1}\text{\rm--}}
\newcommand{\tr}{{\operatorname{tr}}}
\newcommand{\proper}{{\operatorname{prop}}}
\newcommand{\dom}{{\operatorname{dom}}}
\newcommand{\rat}{{\operatorname{rat}}}
\newcommand{\sat}{{\operatorname{sat}}}
\newcommand{\eff}{{\operatorname{eff}}}
\newcommand{\fin}{{\operatorname{fin}}}
\renewcommand{\o}{{\operatorname{o}}}


\newcommand{\op}{{\operatorname{op}}}
\newcommand{\norm}{{\operatorname{norm}}}
\newcommand{\red}{{\operatorname{red}}}
\newcommand{\cont}{{\operatorname{cont}}}
\newcommand{\Zar}{{\operatorname{Zar}}}
\newcommand{\Nis}{{\operatorname{Nis}}}
\newcommand{\et}{{\operatorname{\acute{e}t}}}
\newcommand{\tto}{\dashrightarrow}
\newcommand{\inj}{\hookrightarrow}
\newcommand{\Inj}{\lhook\joinrel\longrightarrow}
\newcommand{\surj}{\rightarrow\!\!\!\!\!\rightarrow}
\newcommand{\Surj}{\relbar\joinrel\surj}
\newcommand{\Res}{\operatorname{Res}}
\newcommand{\Jac}{{\operatorname{Jac}}}
\newcommand{\id}{{\operatorname{Id}}}
\newcommand{\li}{{\operatorname{li}}}
\newcommand{\divi}{{\operatorname{div}}}
\newcommand{\adm}{{\operatorname{adm}}}
\newcommand{\Supp}{{\operatorname{Supp}}}
\newcommand{\pd}{{\partial}}
\newcommand{\pmods}[1]{\; (\operatorname{mod}^*#1)}
\newcommand{\codim}{{\operatorname{codim}}}
\newcommand{\ch}{{\operatorname{ch}}}
\newcommand{\Sym}{{\operatorname{Sym}}}
\newcommand{\Tot}{{\operatorname{Tot}}}
\newcommand{\CH}{{\operatorname{CH}}}
\newcommand{\Mip}{M^\infty_+}

\renewcommand{\lim}{\operatornamewithlimits{\varprojlim}}
\newcommand{\colim}{\operatornamewithlimits{\varinjlim}}
\newcommand{\hocolim}{\operatornamewithlimits{hocolim}}
\newcommand{\ol}{\overline}
\newcommand{\holim}{\operatornamewithlimits{holim}}


\newcommand{\car}{\operatorname{char}}

\renewcommand{\phi}{\varphi}
\renewcommand{\epsilon}{\varepsilon}
\renewcommand{\div}{\operatorname{div}}
\newcommand{\gm}{{\operatorname{gm}}}
\newcommand{\WOQ}{W\underline{\sO}_{\Q}}

\newcommand{\Mt}{\operatorname{\overset{M}{\otimes}}}
\newcommand{\PSTt}{\operatorname{\overset{\PST}{\otimes}}}
\newcommand{\NSTt}{\operatorname{\overset{\NST}{\otimes}}}
\newcommand{\MST}{\operatorname{\mathbf{MST}}}
\newcommand{\MNST}{\operatorname{\mathbf{MNST}}}
\newcommand{\MHI}{\operatorname{\mathbf{MHI}}}
\newcommand{\RSt}{\operatorname{\overset{\RS}{\otimes}}}
\newcommand{\MCor}{\operatorname{\mathbf{MCor}}}
\newcommand{\MP}{\operatorname{\mathbf{MSm}}}
\newcommand{\MPS}{\operatorname{\mathbf{MPS}}}
\newcommand{\MPST}{{\operatorname{\mathbf{MPST}}}}
\newcommand{\CI}{{\operatorname{\mathbf{CI}}}}
\newcommand{\SCRec}{{\operatorname{\mathbf{RSC}}}}
\newcommand{\RSC}{\operatorname{\mathbf{RSC}}}
\newcommand{\Bl}{{\mathbf{Bl}}}
\newcommand{\Cb}{\operatorname{\mathbf{Cube}}}
\newcommand{\ECb}{\operatorname{\mathbf{ECube}}}
\newcommand{\Sets}{\operatorname{\mathbf{Sets}}}
\newcommand{\MT}{\operatorname{\mathbf{MT}}}

\newcommand{\ThZMn}{Th(N_{{Z}}{M},n)}
\newcommand{\ThZMnm}{Th(N_{{Z}}{M},n,m)}
\newcommand{\ThZM}{Th(N_{{Z}}{M},1,0)}

\newcommand{\CItNis}{{\CI^\tau_\Nis}}
\newcommand{\ulaNis}{{\ul{\alpha}_\Nis}}
\newcommand{\omegaCI}{{\omega^{\CI}}}
\newcommand{\RecNis}{{\Rec_\Nis}}


\newcommand{\remind}[1]{{\color{cyan}\large \bf {#1}}}
\newcommand{\bcube}{{\ol{\square}}}
\newcommand{\cube}{\square}

\newcommand{\Zpu}{\Z[1/p]}
\newcommand{\Zputr}{\Z[1/p]_{\text{tr}}}
\newcommand{\lpnbcube}{\Zputr(\bcube^{(l/p^{n})})}
\newcommand{\lpnobcube}{\Zputr(\bcube^{(l/p^{n+1})})}

\newcommand{\CONZ}{C_{0f^{-1}Z}^{N}}
\newcommand{\COMZ}{C_{0Z}^{M}}




\newcommand{\M}{\mathbf{M}}
\newcommand{\Mb}{{\overline{M}}}
\newcommand{\Nb}{{\overline{N}}}
\newcommand{\Lb}{{\overline{L}}}
\newcommand{\Zb}{{\overline{Z}}}
\newcommand{\CbS}{{\overline{C}_S}}
\newcommand{\Cbeta}{{\overline{C}_\eta}}
\newcommand{\Gu}{{\underline{G}}}
\def\indlim#1{\underset{{\underset{#1}{\longrightarrow}}}{\mathrm{lim}}\; }
\def\rmapo#1{\overset{#1}{\longrightarrow}}

\newcommand{\Cu}{\mathcal{C}}
\newcommand{\CuX}{\Cu_\Xb}
\newcommand{\CuMN}{\Cu_{(\Mb,N)}}

\newcommand{\ulM}{\underline{\M}}
\newcommand{\ulMP}{\operatorname{\mathbf{\underline{M}Sm}}}
\newcommand{\ulMNS}{\operatorname{\mathbf{\underline{M}NS}}}
\newcommand{\ulMPS}{\operatorname{\mathbf{\underline{M}PS}}}
\newcommand{\ulMPST}{\operatorname{\mathbf{\underline{M}PST}}}
\newcommand{\ulMNST}{\operatorname{\mathbf{\underline{M}NST}}}
\newcommand{\ulMCor}{\operatorname{\mathbf{\underline{M}Cor}}}
\newcommand{\ulMSm}{\operatorname{\mathbf{\underline{M}Sm}}}
\newcommand{\ulPSCH}{\operatorname{\mathbf{\underline{P}Sch}}}
\newcommand{\ulMSCH}{\operatorname{\mathbf{\underline{M}Sch}}}
\newcommand{\ulPVar}{\operatorname{\mathbf{\underline{P}Var}}}
\newcommand{\ulMVar}{\operatorname{\mathbf{\underline{M}Var}}}
\newcommand{\ulMVarCor}{\operatorname{\mathbf{\underline{M}VarCor}}}
\newcommand{\ulMSmCor}{\operatorname{\mathbf{\underline{M}SmCor}}}
\newcommand{\ulPSm}{\operatorname{\mathbf{\underline{P}Sm}}}



\newcommand{\ulomega}{\underline{\omega}}
\newcommand{\olM}{\overline{\M}}
\newcommand{\olMP}{\operatorname{\mathbf{\overline{M}Sm}}}
\newcommand{\olMPS}{\operatorname{\mathbf{\overline{M}PS}}}
\newcommand{\olMPST}{\operatorname{\mathbf{\overline{M}PST}}}
\newcommand{\olMCor}{\operatorname{\mathbf{\overline{M}Cor}}}
\newcommand{\olomega}{\overline{\omega}}
\newcommand{\oulMP}{\operatorname{\mathbf{\underline{\overline{M}}Sm}}}
\newcommand{\oulMPS}{\operatorname{\mathbf{\underline{\overline{M}}PS}}}
\newcommand{\oulMPST}{\operatorname{\mathbf{\underline{\overline{M}}PST}}}
\newcommand{\oulMCor}{\operatorname{\mathbf{\underline{\overline{M}}Cor}}}
\newcommand{\Comp}{\operatorname{\mathbf{Comp}}}
\newcommand{\CompM}{\Comp(M)}
\newcommand{\ulMDMgm}{\ulMDM_{\gm}^{\eff}}
\newcommand{\ulMDMNis}{{\ulMDM_{\Nis}^{\eff}}}


\newcommand{\MpMnZ}{\ol{M},\Minf + n\ol{Z}}
\newcommand{\MpNnZ}{\ol{N},N^\inf + nf^{-1}\ol{Z}}
\def\bN{\mathbb{N}}
\def\bZ{\mathbb{Z}}
\def\bQ{\mathbb{Q}}
\def\Ztr{\bZ_\tr}
\def\Qtr{\bQ_\tr}
\def\soO{\overline{\sO}}
\def\Image{\mathrm{Image}}
\def\rC#1{rel\sC(#1)}
\def\pb{\overline{p}}
\def\Xd#1{X_{(#1)}}
\def\Xbd#1{\Xb_{(#1)}}
\def\Ud#1{U_{(#1)}}
\def\isom{\overset{\cong}{\longrightarrow}}

\def\Xinf{X^\infty}
\def\Minf{M^{\infty}}

\def\Zinf{Z^\infty}
\newcommand{\MCornc}{{\MCor}^{\text{snc}}}
\newcommand{\MCorncf}{{\MCor}^{\text{snc},\text{fin}}}
\newcommand{\MCorls}{{\MCor}_{\text{ls}}}
\newcommand{\ulMCorls}{{\ulMCor}_{\text{ls}}}
\newcommand{\ulMSmls}{{\ulMSm}_{\text{ls}}}
\newcommand{\ulMSmlsCor}{{\ulMSm}_{\text{ls}}\Cor}

\newcommand{\MpM}{(\ol{M},\Minf)}
\newcommand{\MpZ}{(\ol{Z},\Zinf)}

\newcommand{\til}{\widetilde}
\newcommand{\uPic}{\underline{\Pic}}

\newcommand{\Vb}{\overline{V}}
\newcommand{\qa}{{\operatorname{qa}}}

\newcommand{\MMn}{\mathbb{B}_n}
\newcommand{\MpA}{(\mathbb{A}^1,\emptyset)}
\newcommand{\MpnP}{(\mathbb{P}^1,n\{\infty\})}



\def\tV{\widetilde{V}}
\def\tW{\widetilde{W}}
\def\Vb{\overline{V}}
\def\Wb{\overline{W}}
\def\WbNo{\Wb^{N,o}}
%\def\bP{\mathbb{P}}
\def\bP{\mathbf{P}}
\def\tCXY#1#2{\widetilde{C}_{#1}(\Xb|Y)(#2)}
\def\CXY#1#2{C_{#1}(\Xb|Y)(#2)}
\def\tCX#1#2{\widetilde{C}_{#1}(X)(#2)}
\def\CX#1#2{C_{#1}(X)(#2)}
\def\tCCX#1{\widetilde{C}_{#1}(X)}
\def\tCCXY#1#2{\widetilde{C}_{#1}(\Xb|Y)}
\def\CCXY#1{C_{#1}(\Xb|Y)}
\def\CCX#1{C_{#1}(X)}

\def\CVW#1#2{C_{#1}(\Vb|W)(#2)}
\def\CCVW#1{C_{#1}(\Vb|W)}
\def\CCV#1{C_{#1}(V)}
\def\HS#1#2{H^S_{#1}(#2)}

\newcommand{\hookuparrow}{\mathrel{\rotatebox[origin=c]{90}{$\hookrightarrow$}}}
\newcommand{\hookdownarrow}{\mathrel{\rotatebox[origin=c]{-90}{$\hookrightarrow$}}}



\newcounter{spec}
\newenvironment{thlist}{\begin{list}{\rm{(\roman{spec})}}%
{\usecounter{spec}\labelwidth=20pt\itemindent=0pt\labelsep=10pt}}%
{\end{list}}%


% counters with three numbers,like Theorem 1.2.3
% @Keiho: I changed subsection to section, so for example, Lemma 2.0.1 becomes 2.1.
\newtheorem{lemma}{Lemma}[section]
% counters with two numbers, Theorem 1.2
%\newtheorem{lemma}{Lemma}[section]

\newtheorem{hyp}[lemma]{Hypothesis}



\newtheorem{Th}{Theorem}
\newtheorem{Corr}[Th]{Corollary}
\newtheorem{Rm}[Th]{Remark}

%\swapnumbers
%\newtheorem{lemma}{Lemma}[subsection]
%\newtheorem{lemma}{Lemma}[section]
\newtheorem{thm}[lemma]{Theorem}
\newtheorem{theorem}[lemma]{Theorem}
\newtheorem{defnthm}[lemma]{Definition/Theorem}
\newtheorem{prop}[lemma]{Proposition}
\newtheorem{proposition}[lemma]{Proposition}
\newtheorem{cor}[lemma]{Corollary}
\newtheorem{corollary}[lemma]{Corollary}
\newtheorem{conj}[lemma]{Conjecture}
\theoremstyle{definition}
\newtheorem{situ}[lemma]{Situation}
\newtheorem{defn}[lemma]{Definition}
\newtheorem{constr}[lemma]{Construction}
\newtheorem{definition}[lemma]{Definition}
\newtheorem{nota}[lemma]{Notation}
\newtheorem{note}[lemma]{Note}
\newtheorem{figu}[lemma]{Figure}
\newtheorem{para}[lemma]{}
\theoremstyle{remark}
\newtheorem{qn}[lemma]{Question}
\newtheorem{rk}[lemma]{Remark}
\newtheorem{remark}[lemma]{Remark}
\newtheorem{remarks}[lemma]{Remarks}
\newtheorem{ex}[lemma]{Example}
\newtheorem{exs}[lemma]{Examples}
\newtheorem{example}[lemma]{Example}
\newtheorem{claim}[lemma]{Claim}
\newtheorem{exo}[lemma]{Exercise}
\newtheorem{warning}[lemma]{Warning}
\numberwithin{equation}{section}


\newcommand{\udo}{\underline{\omega}}


\setcounter{tocdepth}{1}
\newtheorem{notation}[lemma]{Notation}


\newcommand{\shuji}[1]{\begin{color}{green}{#1}\end{color}}
\newcommand{\shujirem}[1]{\begin{color}{red}{Shuji:\;#1}\end{color}}

\newcommand{\mf}{\mathfrak}


\begin{document}
\title[Voevodsky motives and motives with modulus]{Voevodsky motives and motives with modulus in positive characteristic}
\author{Keiho Matsumoto}

\begin{abstract}
Let $k$ be a perfect field of characteristic $p>0$. In this paper, without assuming resolution of singularities, we prove that the triangulated category of motives with modulus with rational coefficients is equivalent to Voevodsky's triangulated category of motives with rational coefficients:
\[
\MDM^\eff(k,\Q)\simeq \DM^\eff(k,\Q).
\]
Equivalently, after tensoring with $\Q$, the multiplicities of the modulus become invisible in the category of motives with modulus in positive characteristic.
\end{abstract}

\maketitle

\section{Introduction}

Let $k$ be a perfect field of characteristic $p>0$. In the series of papers \cite{KMSY19a}, \cite{KMSY19b}, and \cite{KMSY22a}, Kahn, Miyazaki, Saito, and Yamazaki constructed the triangulated categories of motives with modulus $\MDM^\eff(k)$ and $\ulMDM^\eff(k)$, which extend Voevodsky's triangulated category of motives $\DM^\eff(k)$. Whereas $\DM^\eff(k)$ is built from smooth varieties, the categories $\MDM^\eff(k)$ and $\ulMDM^\eff(k)$ are built from modulus pairs, that is, pairs consisting of a variety together with an effective Cartier divisor. One of the principal motivations for the introduction of motives with modulus is to provide a motivic framework for cohomology theories that are not $\AA^1$-invariant; see, for instance, \cite{KSY15}. A basic example is coherent cohomology $H^i_{\Zar}(-,\sO)$.

It is well known that the Hom-groups in the triangulated category $\DM^\eff$ are expressed in terms of higher Chow groups; see \cite[Corollary 4.2.6]{V02}. It is therefore natural to expect that the Hom-groups in $\ulMDM^\eff$ should be described by cycle-theoretic invariants with modulus, notably the higher Chow groups with modulus introduced by Binda and Saito \cite{BS18}. We stress that this expectation is only heuristic at present: it does not arise from a direct comparison of the corresponding modulus conditions, which point in opposite directions. For a variety $X$ and effective Cartier divisors $D$ and $D'$ on $X$ with the same support, Miyazaki proved in \cite{M19a} that, after inverting $p$, higher Chow groups with modulus are independent of the multiplicities of the modulus:
\[
\CH^i(X|D,j)\bigl[\tfrac{1}{p}\bigr] \simeq \CH^i(X|D',j)\bigl[\tfrac{1}{p}\bigr].
\]
This suggests that the motives $\ulM(X,D)\bigl[\tfrac{1}{p}\bigr]$ and $\ulM(X,D')\bigl[\tfrac{1}{p}\bigr]$ should be isomorphic in $\ulMDM^\eff(k,\Z[\tfrac{1}{p}])$. The main result of this paper establishes a rational version of this expectation.

\begin{thm}\label{main}
Let $k$ be a perfect field of characteristic $p>0$. Then the functor
\[
\underline{\omega}_{\eff}:\ulMDM^\eff(k,\Q)\to \DM^\eff(k,\Q)
\]
induces an equivalence
\[
\MDM^\eff(k,\Q)\simeq \DM^\eff(k,\Q).
\]
In particular, if $X$ is a normal proper variety and $D,D'$ are effective Cartier divisors on $X$ with the same support such that $X\setminus |D|$ is smooth over $k$, then there is an isomorphism
\[
\ulM(X,D)_\Q \simeq \ulM(X,D')_\Q
\]
in $\MDM^\eff(k,\Q)$.
\end{thm}

Theorem~\ref{main} is essentially reduced to the following statement. Without resolution of singularities, this result seems rather delicate. Our proof rests on Kelly--Miyazaki's work on the modulus qfh topology \cite{KM21}, Bhatt--Snowden's refined alteration theorem \cite{BS17}, and the notion of an interior generically finite $h$-covering introduced in the present paper.

\begin{thm}\label{main3}
The triangulated category $\ulMDM^\eff(k,\Q)$ is compactly generated by $\ulM(M)_\Q$ for log smooth modulus pairs $M \in \ulMSmls$. Similarly, the triangulated category $\MDM^\eff(k,\Q)$ is compactly generated by $\ulM(M)_\Q$ for log smooth modulus pairs $M \in \ulMSmls$ such that $\ol{M}$ is proper.
\end{thm}

We briefly explain the proof of Theorem~\ref{main3}. We follow the notation and conventions of \cite{KM21}. A \textit{modulus pair} $M$ over $k$ is a pair $(\ol{M},\Minf)$, where $\ol{M}$ is a separated scheme of finite type over $k$ and $\Minf$ is an effective Cartier divisor on $\ol{M}$. We write
\[
M^\circ = \ol{M}\setminus |M^\infty|
\]
and call it the \textit{interior} of $M$. Let $\ulPSCH(k)$ be the category whose objects are all modulus pairs over $k$. For two modulus pairs $N$ and $M$, a morphism $N\to M$ is a morphism of schemes $\ol{N}\to \ol{M}$ satisfying the modulus condition that the closed immersion
\[
M^\infty \times_{\ol{M}} \ol{N} \to \ol{N}
\]
factors through $N^\infty$. We call a morphism $f:N \to M$ \textit{minimal} if the natural morphism
\[
M^\infty \times_{\ol{M}} \ol{N} \to N^\infty
\]
is an isomorphism. We define the category of \textit{smooth modulus pairs} $\ulPSm$ to be the full subcategory of $\ulPSCH$ consisting of those modulus pairs $(\ol{M},\Minf)$ for which the interior $M^\circ$ is smooth over $k$. Let $\ul{\Sigma}$ be the class of minimal morphisms $s:X' \to X$ in $\ulPSCH$ such that $\ol{X'} \to \ol{X}$ is proper and $s^{-1}(X^\circ)\simeq X^\circ$. We then set
\[
\ulMSCH:= \ulPSCH[(\ul{\Sigma})^{-1}], \qquad \ulMSm:= \ulPSm[(\ul{\Sigma})^{-1}].
\]
A \textit{log smooth modulus pair} is a modulus pair $(\ol{M},M^\infty)$ such that $\ol{M}$ is smooth over $k$ and the support $|M^\infty|$ is a strict normal crossings divisor on $\ol{M}$. We write $\ulMSmls$ for the full subcategory of $\ulMSm$ spanned by such pairs.

In Section~\ref{section2}, we introduce the notion of an \textit{interior generically finite $h$-covering} on $\ulMSCH$, which is stronger than the modulus qfh covering of Kelly--Miyazaki \cite[Definition 4.5]{KM21}. This notion has two essential features. First, it allows one to cover a smooth modulus pair by log smooth modulus pairs. More precisely, we prove that for every smooth modulus pair $M$, there exists an interior generically finite $h$-covering
\[
\{N_i\to M \}_{i \in I}
\]
such that each $N_i$ is a log smooth modulus pair; see Proposition~\ref{altalt}. The proof uses Bhatt--Snowden's refined alteration theorem \cite{BS17}. Secondly, up to $\ul{\Sigma}$, every interior generically finite $h$-covering can be refined by a modulus qfh covering; see Proposition~\ref{hisqfh}.

In Section~\ref{section3}, combining interior generically finite $h$-coverings with the equivalences of Kelly--Miyazaki \cite[Theorems 6.9 and 6.10]{KM21},
\begin{equation}\label{intro3}
\Shv_{\Mqfh}(\ulMSCH(k),\Q) \simeq \Shv_{\Mqfh}(\ulMCor(k),\Q) \simeq \Shv_{\ulM \Nis}(\ulMCor(k),\Q),
\end{equation}
we prove that for every smooth modulus pair $M \in \ulMSm$, there exists a finite family of log smooth modulus pairs $\{N_i\}_{i \in I}$ such that $\Q_\tr(M)$ is a direct summand of $\bigoplus_i \Q_\tr(N_i)$ in $\Shv_{\ulM\Nis}(\ulMSmCor,\Q)$. Theorem~\ref{main3} follows formally from this.

\section*{Acknowledgements}
We thank Shane Kelly for helpful discussions on this circle of ideas, for his earlier suggestion in 2019 of a notion closely related to the interior generically finite $h$-site on $\ulMCor$, and for helpful comments on a previous version of this paper. We also thank Hiroyasu Miyazaki for helpful comments on an earlier draft.

\section*{Notations}

For a modulus pair $M$, we denote by $\ol{M}$ its total space, by $M^\infty$ its modulus, and by $M^\circ=\ol{M} \setminus |M^\infty|$ its interior. For a small category $\sC$ and a Grothendieck topology $\tau$ on $\sC$, we write $\Shv_{\tau}(\sC)$ for the category of $\tau$-sheaves on $\sC$ with values in $\Z\text{-Mod}$, and $\Shv_{\tau}(\sC,\Q)$ for the category of $\tau$-sheaves on $\sC$ with values in $\Q\text{-Vect}$.

\section{The interior generically finite \texorpdfstring{$h$}{h}-site}\label{section2}

We introduce a class of coverings on $\ulMSCH$ that is stronger than the modulus qfh topology of Kelly--Miyazaki \cite[Definition 4.7]{KM21}.

Let $\mathcal{U}$ be a finite family of minimal morphisms in $\ulPSCH$. We say that $\mathcal{U}$ is a \textit{proper interior generically finite covering} if it is of the form
\[
\{(\ol{U}_{i},U_{i}^\infty) \overset{f_i}{\to} (\ol{M},\Minf)\}_{i \in I},
\]
where $I$ is finite, each $\ol{f}_i:\ol{U}_i \to \ol{M}$ is proper and generically finite, the induced morphism
\[
\coprod_{i\in I}\ol{U}_i \longrightarrow \ol{M}
\]
is surjective, and the following condition is satisfied:
\begin{itemize}
    \item[($\star$)] For every point $x\in \ol{M}\setminus |\Minf|$, there exist an open neighbourhood $V \subset \ol{M}\setminus |\Minf|$ of $x$ and an index $i\in I$ such that the base change
    \[
    \ol{U}_{i}\times_{\ol{M}} V \to V
    \]
    is finite surjective.
\end{itemize}

We say that a family $\mathcal{U}$ is an \textit{interior generically finite $h$-covering} if it is of the form
\[
\{U_{i_0i_1,\dots,i_n} \to \cdots \to U_{i_0i_1} \to U_{i_0}\}_{i_0 \in I_0,\, i_1\in I_{i_0},\,\dots,\,i_n\in I_{i_{n-1}}},
\]
where at each stage the corresponding family is either a Zariski open covering or a proper interior generically finite covering.

The class of interior generically finite $h$-coverings in $\ulMSCH$ is defined to be the smallest class of families of morphisms containing the images of interior generically finite $h$-coverings in $\ulPSCH$ and stable under isomorphisms and composition.

We first show that a single proper interior generically finite covering in $\ulPSCH$ may be replaced, up to morphisms in $\ul{\Sigma}$, by a finite covering.

\begin{lemma}\label{singlelemma}
For any single proper interior generically finite covering $\{f:N \to M\}$ in $\ulPSCH$, there exist a modulus pair $V$, a finite covering $\{V\to M\}$, and a morphism $r:N \to V$ in $\ul{\Sigma}$ fitting into the commutative diagram
\[
\xymatrix{
N \ar[d]_-{r \in \ul{\Sigma}} \ar[drr]^-{f}& & \\
V \ar[rr]_-{\text{finite covering}}&  & M.
}
\]
\end{lemma}

\begin{proof}
Since $f:N \to M$ is a single proper interior generically finite covering, the morphism $\ol{f}:\ol{N}\to \ol{M}$ is proper surjective and its restriction
\[
f^\circ:N^\circ \to M^\circ
\]
is finite surjective. Let
\[
\ol{{V}} := \underline{\Spec}_{\sO_{\ol{M}}}\ol{f}_*\sO_{\ol{N}},
\]
and let $g:\ol{V}\to \ol{M}$ be the induced finite morphism. Endow $\ol{V}$ with the divisor $g^*\Minf$. Then
\[
(\ol{V},g^*\Minf) \to (\ol{M},\Minf)
\]
is a finite covering. By the universal property of relative $\underline{\Spec}$, the morphism $\ol{f}$ factors canonically through a proper morphism
\[
\ol{r}:\ol{N}\to \ol{V},
\]
and this induces a morphism of modulus pairs
\[
r:N \to (\ol{V},g^*\Minf).
\]
Since $f$ is minimal, so is $r$. Moreover, over the interior we have
\[
V^\circ=\underline{\Spec}_{\sO_{M^\circ}}(f^\circ_*\sO_{N^\circ}),
\]
and because $f^\circ$ is finite, the canonical morphism $N^\circ\to V^\circ$ is an isomorphism. Therefore $r$ belongs to $\ul{\Sigma}$.
\end{proof}

We next show that every interior generically finite $h$-covering can be refined, up to $\ul{\Sigma}$, by a modulus qfh covering.

\begin{prop}\label{hisqfh}
For any proper interior generically finite covering $\{N_i \to M\}_{i\in I}$ in $\ulPSCH$, there exist a modulus pair $M'$, a morphism $b:M' \to M$ in $\ulPSCH$ belonging to $\ul{\Sigma}$, a finite set $J$, an inclusion $J \overset{i_{-}}{\hookrightarrow} I$, an open covering $\{U_j' \to M'\}_{j\in J}$, and for each $j \in J$ a finite covering $\{V_j \to U_j'\}$, a modulus pair $W_j$, a morphism of modulus pairs $r_j: W_j \to V_j$ in $\ulPSCH$, and a morphism $p_j:W_j \to N_{i_j}$ such that:
\begin{itemize}
    \item[(1)] each $r_j$ belongs to $\ul{\Sigma}$;
    \item[(2)] the family
    \[
    \{W_j \overset{r_j}{\to} V_j \to U_j' \to M' \}_{j\in J}
    \]
    is a interior generically finite $h$-covering in $\ulPSCH$, and it refines $\{N_i \to M\}_{i\in I}$ in $\ulMSCH$.
\end{itemize}
In particular, these morphisms fit into the commutative diagram
\[
\xymatrix{
\{W_j\}_{j\in J}\ar[dd]^-{\text{refinement in }\ulMSCH} \ar[rr]^-{r_j \in \ul{\Sigma}}&& \{V_j\}_{j\in J}\ar[rr]^-{\text{finite covering}} && \{U_j'\}_{j\in J}\ar[d]^-{\text{Zariski covering}} \\
&&&&M' \ar[d]^-{b\in \ul{\Sigma}}\\
\{N_i\}_{i\in I}\ar[rrrr]_-{\text{proper interior generically finite covering}}&& &&M.
}
\]
\end{prop}

\begin{proof}
For each point $y \in  M^\circ$, there exist an index ${i_y}\in I$ and an open neighbourhood $y \in U_y \subset  M^\circ$ such that the base change of $f_{i_y}:\ol{N}_{i_y} \to \ol{M}$ to $U_y$ is finite. These open subsets form an open covering
\[
\bigcup_{y \in M^\circ} U_y=\ol{M}\setminus |M^\infty|.
\]
Since $\ol{M}\setminus |M^\infty|$ is of finite type over $k$, we may choose a finite subset $J \subset  M^\circ$ such that
\[
\bigcup_{y \in J} U_y= M^\circ.
\]
We may moreover assume that $U_{y_1}\not\subset U_{y_2}$ for any distinct $y_1,y_2\in J$.

We introduce the notation
\begin{eqnarray*}
Z_y&:=&\text{the closure of }M^\circ \setminus U_y \text{ in }\ol{M},\\
W&:=&\bigcap_{y \in J} Z_y.
\end{eqnarray*}
Since $\bigcup_{y\in J}U_y=M^\circ$, we have $W \cap M^\circ=\emptyset$. Let $\eta_1,\eta_2,\dots,\eta_n$ be the generic points of $|\Minf|$. Each $M^\circ\setminus U_y$ avoids all the $\eta_i$ as well as the generic point of $\ol{M}$. Hence $W$ is properly contained in $|M^\infty|$, and in particular
\[
\codim_{\ol{M}}W \geq 2.
\]

Consider the blowing up
\[
b:\ol{M'}=Bl_{W}\ol{M} \to \ol{M}.
\]
We use the notation
\begin{eqnarray*}
{M'}^\infty&:=& b^*M^\infty,\\
\widetilde{Z_y}&:=&\text{the strict transform of }Z_y\text{ along }b,\\
\ol{U_y'}&:=&\ol{M'}\setminus |\widetilde{Z_y}|,\\
{U'_y}^\infty&:=&{M'}^\infty \cap \ol{U_y'}.
\end{eqnarray*}
By construction,
\[
\bigcap_{y\in J}\widetilde{Z_y}=\emptyset,
\]
so the family $\{\ol{U_y'}\}_{y\in J}$ covers $\ol{M'}$. We also note that $b$ induces an isomorphism
\[
\ol{U_y'}\setminus |{U_y'}^\infty| \simeq U_y.
\]

For $y\in J$, let $U_y'$ denote the modulus pair $(\ol{U_y'},{U'_y}^\infty)$. Then $\{U_y' \to M' \}_{y \in J}$ is a Zariski covering. Let $C_1,C_2,\dots,C_m$ be the connected components of $\ol{U}'_y\times_{\ol{M}}\ol{N}_{i_y}$. Set
\[
\ol{W_y}=\coprod_{i=1}^m (C_i)_{\red}.
\]
Let $p_{y}$ and $g_y$ denote the natural morphisms
\[
\ol{W}_y \to \ol{N}_{i_y}, \qquad \ol{W}_y \to \ol{U}'_{y},
\]
respectively. We then have a diagram
\[
\xymatrix{
\ol{W_y}\ar[rrdd]_{p_y}\ar[rr] \ar@/^20pt/[rrr]^{g_y} && \ol{N}_{i_y} \times_{\ol{M}} \ol{U}_y'  \ar[dd] \ar[r]&\ol{U}_y'\ar@{^{(}-_>}[d]^{\text{open immersion}}\\
&&&\ol{M}' \ar[d]^{b}\\
&& \ol{N}_{i_y} \ar[r]_{f_{i_y}} & \ol{M}.
}
\]
Put
\[
W_y^\infty = g_y^* {U_y'}^\infty.
\]
Since $g_y$ is proper and generically finite, and since $f_{i_y}$ is finite surjective over $U_y$, the morphism $g_y$ is finite surjective over $\ol{U_y'}\setminus |{U_y'}^\infty|$. Therefore
\[
\{(\ol{W}_{y},W_y^\infty) \to (\ol{U_y'},{U_y'}^\infty )\}
\]
is a single proper interior generically finite covering. By Lemma~\ref{singlelemma}, there exist a modulus pair $V_y$, a finite covering
\[
\{V_y \to (\ol{U_y'},{U_y'}^\infty )\},
\]
and a morphism
\[
r_y:(\ol{W}_{y},W_y^\infty) \to V_y
\]
belonging to $\ul{\Sigma}$ such that
\[
\xymatrix{
\ol{W_y}\ar[rrdd]_{p_y}\ar[rr]^-{{r_y \in \ul{\Sigma}}} \ar@/^20pt/[rrr]^{g_y} && V_y  \ar[r]_-{\text{finite covering}} &\ol{U}_y'\ar@{^{(}-_>}[d]^{\text{open immersion}}\\
&&&\ol{M}' \ar[d]^{b}\\
&& \ol{N}_{i_y} \ar[r]_{f_{i_y}} & \ol{M}
}
\]
commutes.

Since $\{(\ol{U_y'}, {U_y'}^\infty) \to (\ol{M'},M'^\infty)\}_{y\in J}$ is a Zariski covering and $b$ belongs to $\ul{\Sigma}$, this is precisely the required refinement.
\end{proof}

\section{Refined alteration and modulus pairs}\label{section3}

In this section, we prove that the derived category $D(\Shv_{\ulM \Nis}(\ulMSmCor,\Q))$ is compactly generated by $\Q_\tr(M)$ for log smooth modulus pairs $M \in \ulMSmls$. Here $\Q_\tr(M)$ denotes the presheaf with transfers represented by $M$ and tensored with $\Q$. The geometric input is Bhatt--Snowden's refined alteration theorem.

\begin{thm}[Refined alteration {\cite[Theorem 2.5]{BS17}}]\label{Bhatt-Snowden}
Let $X$ be a variety over a perfect field $k$. Fix a closed subvariety $Z\hookrightarrow X$ and let $F\subset (X\setminus |Z|)^{\mathrm{sm}}$ be a finite subset of closed points of the smooth locus. Then there exists an alteration $f:Y\to X$ such that $Y$ is smooth, $f^{-1}(Z)$ is a strict normal crossings divisor on $Y$, and $f$ is finite \'etale in a neighbourhood of $F$.
\end{thm}

\begin{prop}\label{altalt}
For every smooth modulus pair $M \in \ulMSm$, there exists a proper interior generically finite covering $\{N_i \to M\}_{i\in I}$ such that each $N_i$ is a log smooth modulus pair.
\end{prop}

\begin{proof}
Let
\[
\ol{X}=\coprod (C_i)_{\red}
\]
be the disjoint union of the reduced irreducible components $C_1,C_2,\dots,C_m$ of $\ol{M}$. Let $p:\ol{X} \to \ol{M}$ be the natural morphism, put $X^\infty=p^*M^\infty$, and write
\[
X=(\ol{X},p^*M^\infty).
\]
It is straightforward to check that the morphism $X\to M$ is a proper interior generically finite covering. Thus it suffices to prove the assertion for $X$. We may assume that $\ol{X}$ is connected. Observe that the interior
\[
X^\circ=\ol{X}\setminus |X^\infty|
\]
is smooth over $k$.

For each point $y \in X^\circ$, Theorem~\ref{Bhatt-Snowden} yields an open neighbourhood $y \in U_y \subset  X^\circ$ and an alteration $f_y:\ol{N}_y \to \ol{X}$ such that $\ol{N}_y$ is smooth, $f_y^{-1}|X^\infty|$ is a strict normal crossings divisor on $\ol{N}_y$, and $f_y$ is finite over $U_y$. Since the opens $U_y$ cover $X^\circ$ and $X^\circ$ is of finite type over $k$, we may choose a finite subset $J \subset  X^\circ$ such that
\[
\bigcup_{y \in J} U_y= X^\circ.
\]
Then the family
\[
\{(\ol{N}_y,f_y^*X^\infty) \to X\}_{y \in J}
\]
is a proper interior generically finite covering, and each term is a log smooth modulus pair.
\end{proof}

In \cite[Corollary 5.18]{KM21}, Kelly--Miyazaki define a monoidal category $\ulMCor$, which may be viewed as a version of $\ulMSCH$ with transfers, together with the categories of modulus qfh sheaves $\Shv_{\Mqfh}(\ulMCor)$ and modulus Nisnevich sheaves $\Shv_{\ulM \Nis}(\ulMCor)$. Let $\ulMSmCor \hookrightarrow \ulMCor$ be the full subcategory consisting of modulus pairs over $k$ with integral normal total space and smooth interior. In \cite{KMSY19a}, Kahn--Miyazaki--Saito--Yamazaki introduce the category of smooth modulus pairs $\ulMCor^{\mathrm{KMSY}}$ over $k$. To avoid confusion, we write $\ulMCor^{\mathrm{KMSY}}$ for their category, although in \cite{KMSY19a} it is denoted by $\ulMCor$. Note that $\ulMSmCor$ is equivalent to $\ulMCor^{\mathrm{KMSY}}$; see \cite{KMSY19a} and \cite[Theorem 5.35]{KM21}.

\begin{prop}\label{kawarikawari}
For every smooth modulus pair $M\in \ulMSmCor$, there exist log smooth modulus pairs $N_1,N_2,\dots,N_n$ such that $\Q_\tr(M)$ is a direct summand of $\Q_\tr(\coprod N_i)$ in $\Shv_{\ulM\Nis}(\ulMSmCor,\Q)$.
\end{prop}

\begin{remark}
The presheaf $\Q_\tr(M)$ is already a modulus Nisnevich sheaf; see \cite[Corollary 5.34]{KM21}. In particular, no further sheafification is needed in the statement of Proposition~\ref{kawarikawari}.
\end{remark}

\begin{proof}
Since the functor $\ulMSmCor \to\ulMCor$ is fully faithful, the induced left Kan extension defines a fully faithful functor
\[
\Shv_{\ulM \Nis}(\ulMSmCor, \Q) \to \Shv_{\ulM \Nis}(\ulMCor, \Q).
\]
By \cite[Theorems 6.9 and 6.10]{KM21}, there are equivalences
\[
\Shv_{\ulM\Nis}(\ulMCor,\Q)\simeq \Shv_{\Mqfh}(\ulMCor,\Q) \simeq \Shv_{\Mqfh}(\ulMSCH,\Q),
\]
under which $\Q_\tr(M)$ is sent to $\Q_{\Mqfh}(M)$, the modulus qfh sheafification of $\Q(M)$. It therefore suffices to show that for every smooth modulus pair $M$, there exist log smooth modulus pairs $N_1,N_2,\dots,N_n$ such that $\Q_{\Mqfh}(M)$ is a direct summand of $\Q_{\Mqfh}(\coprod N_i)$ in $\Shv_{\Mqfh}(\ulMCor,\Q)$.

Let $M\in \ulMSm$ be smooth. By Proposition~\ref{altalt}, there exists a proper interior generically finite covering $\{N_i \to M\}_{i\in I}$ such that each $N_i$ is log smooth. By Proposition~\ref{hisqfh}, after replacing this family by a refinement, we may find a subset $J \overset{i_{-}}{\hookrightarrow} I$, a modulus qfh covering $\{W_j \to M\}_{j \in J}$ in $\ulMSCH$, and proper interior generically finite coverings $W_j \to N_{i_j}$ for $j\in J$, satisfying the conditions of Proposition~\ref{hisqfh}. The morphism $\coprod_{j\in J} W_j \to M$ induces a natural map
\[
\Hom_{\Shv_{\Mqfh}(\ulMSCH,\Q)}(\Q_{\Mqfh}(M),\Q_{\Mqfh}(\coprod_{j\in J} W_j)) \to \Hom_{\Shv_{\Mqfh}(\ulMSCH,\Q)}(\Q_{\Mqfh}(M),\Q_{\Mqfh}(M)),
\]
which is surjective, because $\Q_{\Mqfh}(\coprod_{j\in J} W_j) \to \Q_{\Mqfh}(M)$ is the sheafification of the covering morphism $\coprod_{j\in J} W_j \to M$; see also \cite[Proof of Theorem 6.10]{KM21}. Since the morphism $\coprod_{j\in J} W_j \to M$ factors through $\coprod_{i\in I} N_i \to M$, the map
\[
\Hom_{\Shv_{\Mqfh}(\ulMSCH,\Q)}(\Q_{\Mqfh}(M),\Q_{\Mqfh}(\coprod_{i\in I} N_i)) \to \Hom_{\Shv_{\Mqfh}(\ulMSCH,\Q)}(\Q_{\Mqfh}(M),\Q_{\Mqfh}(M))
\]
is likewise surjective. Hence the surjection
\[
\Q_{\Mqfh}(\coprod_{i\in I} N_i) \to \Q_{\Mqfh}(M)
\]
admits a section, so $\Q_{\Mqfh}(M)$ is a direct summand of $\Q_{\Mqfh}(\coprod_{i\in I} N_i)$. This proves the assertion.
\end{proof}

By the same argument, one obtains the following proper version.

\begin{prop}\label{joker}
For every smooth modulus pair $M\in \ulMSmCor$ such that $\ol{M}$ is proper, there exist log smooth modulus pairs $N_1,N_2,\dots,N_n$ such that each $\ol{N}_i$ is proper and $\Q_\tr(M)$ is a direct summand of $\Q_\tr(\coprod N_i)$ in $\Shv_{\ulM\Nis}(\ulMSmCor,\Q)$.
\end{prop}

We briefly recall the definition of $\ulMDM^\eff$ from \cite{KMSY22a}. The triangulated category $\ulMDM^\eff$ is the Verdier quotient of $D(\Shv_{\ulM \Nis}(\ulMSmCor))$ by the smallest localizing subcategory containing all complexes of the form

(CI) for ${\mathfrak{M}}\in\ulMCor$,
\[
\Ztr(\mathfrak{M}\otimes \cube) \to \Ztr(\mathfrak{M}).
\]
We define $\MDM^\eff$ to be the smallest full triangulated subcategory of $\ulMDM^\eff$ containing the objects $\ulM(\ol{M},\Minf)$ for modulus pairs $(\ol{M}, \Minf)$ with $\ol{M}$ proper, and stable under isomorphisms, direct sums, shifts, and cones. In \cite{KMSY22a}, Kahn--Miyazaki--Saito--Yamazaki construct $\MDM^\eff$ in a different, more direct fashion. The two definitions are equivalent by \cite[Theorem 5.2.2]{KMSY22a}.

\begin{thm}\label{lsequalsch}
The triangulated category $\ulMDM^\eff(k,\Q)$ is compactly generated by $\ulM(M)_\Q$ for log smooth modulus pairs $M \in \ulMSmls$. Similarly, the triangulated category $\MDM^\eff(k,\Q)$ is compactly generated by $\ulM(M)_\Q$ for log smooth modulus pairs $M \in \ulMSmls$ such that $\ol{M}$ is proper.
\end{thm}

\begin{proof}
By \cite[Theorem 4.1.1]{KMSY22a}, the triangulated category $\ulMDM^\eff(k,\Q)$ is compactly generated by $\ulM(M)_\Q$ for smooth modulus pairs $M\in \ulMSm$. Likewise, $\MDM^\eff(k,\Q)$ is compactly generated by $\ulM(M)_\Q$ for smooth modulus pairs $M\in \ulMSm$ such that $\ol{M}$ is proper. The claim therefore follows from Proposition~\ref{kawarikawari} and Proposition~\ref{joker}.
\end{proof}


There is an adjunction $\ul{\omega}_{\eff}:\ulMDM^\eff \rightleftarrows\DM^{\eff}:\ul{\omega}^\eff$, where $\ul{\omega}_{\eff}$ is a localization functor and $\ul{\omega}^\eff$ is fully faithful (see \cite[Section 6]{KMSY22a}). In our previous paper \cite[Theorem~7.1]{M22}, we prove that for a smooth variety $M^\circ$ over $k$ which has a compactification $\ol{M}$ such that $\ol{M}$ is smooth and $|\ol{M}\backslash M^\circ|$ is a strict normal crossings divisor on $\ol{M}$, the unit
\begin{equation}\label{unit}
\ulM(\ol{M},|\ol{M}\backslash M|_{\red}) \to \ul{\omega}^{\eff}({\bf{M}}(M^\circ))
\end{equation}
of the adjunction is an isomorphism. Notice that the functor 
\[
\omega_\eff:\MDM^\eff \hookrightarrow \ulMDM^\eff \overset{\ul{\omega}_\eff}{\to} \DM^\eff
\]
admits the right-adjoint functor $\omega^\eff$, which is strongly additive and fully faithful \cite[Proposition 6.1.2]{KMSY22a}.

    \begin{thm}[Comparison Theorem]\label{mainmainmain}
    Let $k$ be a perfect field of positive characteristic. Then the localization functor
    \[
    \omega_\eff:\MDM^\eff(k,\Q) \to \DM^\eff(k,\Q)
    \]
    is an equivalence.
\end{thm}
\begin{proof}
By Theorem~\ref{lsequalsch}, $\MDM^\eff(k,\Q)$ is compactly generated by $\ulM(M)_{\Q}$ for proper log smooth modulus pairs $M=(\ol{M},\Minf) \in \ulMSmls$. By \cite[Corollary 8.8]{M22}, $\ulM(\ol{M},\Minf)_\Q \simeq (\ol{M},|\Minf|)_\Q$ in $\ulMDM^\eff(k,\Q)$. The functor $\omega^\eff$ is strongly additive and fully faithful, so it is enough to show that for every proper log smooth modulus pair $M=(\ol{M},M^\infty)$ with $M^\infty$ reduced, the object $\ulM(\ol{M},M^\infty)_{\Q}$ belongs to the essential image of $\omega^\eff$. For such $M$, the isomorphism \eqref{unit} shows that
\[
\ulM(\ol{M},M^\infty)_{\Q}\simeq \omega^\eff(\mathbf{M}(M^\circ)_{\Q}).
\]
Hence all reduced proper log smooth generators lie in the essential image of $\omega^\eff$, which proves the claim.\end{proof}
\begin{cor}\label{corcorcor}
     For a normal proper variety $X$ and effective Cartier divisors $D$ and $D'$ such that the supports
$|D|$ and $|D'|$ coincide, and $X\backslash |D|$ is smooth over $k$, there is an isomorphism in $\MDM^\eff(k,\Q)$:
\[
\ulM(X,D)_{\Q} \simeq \ulM(X,D')_\Q.
\]
\end{cor}
\begin{proof}
    The functor $\omega_\eff$ is an equivalence and sends $\ulM(X,D)_{\Q}$ and $\ulM(X,D')_{\Q}$ to the same target $\textbf{M}(X\backslash |D|)_{\Q}$ (see \cite[Proposition 6.1.1]{KMSY22a}).
\end{proof}


\bibliography{bib}
\bibliographystyle{plain}

\end{document}
